% ---------------------------------------------------------------------------
% Results section for the VizWiz-Captions captioning study.
% References the experimental design of Section 4.5 and the tables in
% paper_results_tables.tex. Uses \ref{...}; if you load cleveref, you may
% replace "Table~\ref{...}" with "\Cref{...}".
% ---------------------------------------------------------------------------

\section{Results}
\label{sec:results}

We present results following the four axes of the experimental design
(Section~\ref{sec:design}): \emph{where} to adapt the backbone
(Section~\ref{sec:res-scaling}), whether OCR signal helps
(Section~\ref{sec:res-ocr}), whether OCR and vision adaptation combine
(Section~\ref{sec:res-ade-vision}), and whether BLIP-base is the right backbone
(Section~\ref{sec:res-backbones}). All numbers are CIDEr on VizWiz
\texttt{val} unless noted, and every comparison in Sections~\ref{sec:res-ocr}
and~\ref{sec:res-ade-vision} is a controlled A/B in which only the stated factor
varies.

\subsection{Where to adapt: scaling the LoRA target}
\label{sec:res-scaling}

Table~\ref{tab:main-results} reports the main fine-tuning ladder. Parameter-%
efficient fine-tuning already lifts the un-adapted BLIP-base from
43.5 to 60.3 CIDEr ($+16.8$) with text-only LoRA on the full training set, of
which $+5.9$ comes from scaling the data from 20k pairs to the full corpus
($54.4\!\to\!60.3$). The decisive lever, however, is \emph{where} adaptation is
applied rather than how much data it sees: extending LoRA to the vision
encoder's attention projections (adding \texttt{qkv}, rank $r{=}32$) raises CIDEr
from 60.3 to 72.7, a $+12.4$ jump that dominates every other factor we study.
Adapting the visual features to the domain --- blurry, oddly framed, poorly lit
photographs --- matters far more than further adapting the text decoder.

Beyond this point the configuration saturates. Extending training from 2 to 4
epochs adds a further $+1.9$ (to \textbf{74.6} CIDEr) with no sign of
overfitting, while raising the rank to $r{=}64$ does not help (74.0 vs.\ 74.6),
confirming $r{=}32$ as the operating point. A learning-rate sweep reproduces the
need to validate at the final budget: $\text{lr}{=}2\mathrm{e}{-}4$ wins at 2
epochs (73.5 vs.\ 72.7) but \emph{loses} at 4 epochs (73.7 vs.\ 74.6), so the
ranking inverts with budget. The resulting model --- text+vision LoRA,
$r{=}32$, 4 epochs, $\text{lr}{=}1\mathrm{e}{-}4$, dropout 0.05 --- is our
\textbf{best autonomous model} (74.6 CIDEr, free at inference) and serves as the
control for the OCR experiments that follow.

\subsection{Does OCR signal help? (Phases 1--4b)}
\label{sec:res-ocr}

\paragraph{Engine characterization (Phase~1).}
Since VizWiz-Captions carries no OCR ground truth, we first characterize the two
engines qualitatively (Table~\ref{tab:ocr-engines}). EasyOCR fails or returns
empty on 6 of 30 with-text images and collapses into noise on degraded photos,
whereas Landing AI ADE never fails and extracts $3.4\times$ more text. Crucially,
ADE is not a pure OCR engine: on 20 of 30 images it returns a \emph{scene
description} rather than a verbatim transcription, placing its output closer to a
caption than to raw text --- a distinction that turns out to drive every
subsequent result.

\paragraph{A literal OCR prefix hurts (Phase~2).}
Injecting EasyOCR text as a conditional decoder prefix \emph{degrades}
captioning (Table~\ref{tab:phase2-easyocr}): $-3.35$ CIDEr globally and $-4.31$
on the with-text subset. The experiment is clean --- the no-text subset, where
no prefix is injected, is unchanged ($+0.25$) --- so the damage comes from the
noisy OCR \emph{content}, not from the prefixing mechanism.

\paragraph{Quality is decisive, and it is the description that helps (Phase~3).}
Replacing EasyOCR with ADE produces a strictly monotone ordering
(Table~\ref{tab:phase3-ade}): EasyOCR (noisy) $<$ ADE text-only (clean but
literal) $<$ full ADE, and only full ADE beats the control ($+2.5$ globally,
$+3.2$ on with-text). Because the literal, text-only variant remains \emph{below}
the control, the gain is not attributable to ``reading the text'': it comes from
ADE's scene description, which effectively distills a strong external
vision--language model. The winning variant, however, carries a deployment cost
--- it requires an ADE call per image at inference.

\paragraph{The gain does not distill into an autonomous model (Phases~4, 4b).}
We therefore attempt to internalize the Phase-3 gain without the inference-time
dependency. Using ADE descriptions directly as pseudo-caption targets
(Phase~4) is sharply counterproductive (Table~\ref{tab:phase4-distill}): CIDEr
collapses from 55.1 to 33.5 as the student absorbs ADE's verbose prose and
applies it everywhere, degrading even the no-text subset
($62.9\!\to\!46.7$) where no pseudo-captions were added. Correcting the
caption-style mismatch --- distilling from a teacher that already emits short,
VizWiz-style captions (Phase~4b) --- recovers the collapse
($33.5\!\to\!54.2$, Table~\ref{tab:phase4b-distill}) but only reaches parity with
the control and stays below Phase~3. In short, style matching neutralizes the
damage without capturing the gain: the value of the strong VLM resides in its
per-image inference, which a small student does not internalize from
$\sim\!2{,}000$ pseudo-labels.

\subsection{Combining ADE with vision adaptation}
\label{sec:res-ade-vision}

The Phase-3 gain was measured on a \emph{weak} backbone (20k, text-only,
$r{=}16$), leaving open whether ADE is still useful once the vision encoder is
adapted, or whether vision scaling already captures the same object/brand
information. Phase~A answers this at the 20k scale
(Table~\ref{tab:phaseA-ade-vision}): the signals are additive. ADE adds $+4.52$
CIDEr on with-text images \emph{over} the vision-adapted control --- larger than
the $+3.2$ seen on the text-only backbone --- so vision scaling does not subsume
ADE. EasyOCR still hurts ($-3.04$), though less than on the text-only backbone,
confirming that better visual features ignore noise more gracefully without
making it useful.

Phase~B confirms the effect at full-train scale
(Table~\ref{tab:phaseB-ade-vision}): the $\sim\!+4$ delta holds, yielding our
\textbf{best overall result of 78.68 CIDEr} ($+4.08$ over the 74.6 vision-only
model). A near-perfect symmetry emerges around the control --- noisy OCR
subtracts ${\sim}4.3$, strong OCR adds ${\sim}4.1$ --- and the ordering
EasyOCR $<$ control $<$ ADE replicates in both regimes. The practical trade-off
is explicit: the 78.68 model requires a paid ADE call per image at inference,
so we report two operating points --- \textbf{78.68} (best overall, paid) and
\textbf{74.60} (best autonomous, free and local).

\subsection{Alternative vision--language backbones}
\label{sec:res-backbones}

Finally, we verify that BLIP-base is a reasonable backbone choice
(Table~\ref{tab:alt-vlm}). GIT-large is more data-efficient --- it outperforms
BLIP at 20k (69.3 vs.\ 64.7) --- but this advantage does not survive scaling:
with rank matched at $r{=}32$ on the full set, GIT reaches 73.9 against BLIP's
74.6, a near-tie with BLIP marginally ahead. This is a cautionary case for
low-data prototyping: the 20k ranking would have led to the wrong conclusion.
Qwen2.5-VL-3B is feasible in 8\,GB via 4-bit QLoRA but $\sim\!14\times$ slower,
and its bilingual prior leaks Chinese into 26\% of captions; blocking CJK tokens
at generation raises it to 63.6 CIDEr (64.9 on with-text, 55.8 on no-text),
better than BLIP on with-text but worse without, and still below both BLIP and
GIT at 20k --- and forcing English in the \emph{training} prompt (an 11-hour
rerun) made it worse, not better. Florence-2-large could not be evaluated: its
custom remote code is incompatible with the installed Transformers release. We
conclude that no alternative dethrones BLIP-base for this task within the 8\,GB
budget, and that BLIP-base was the correct backbone choice.
